\documentclass[a4paper]{article}

\usepackage[spanish]{babel}
\usepackage{graphicx}
\usepackage{amsmath, amssymb}
\usepackage[margin=2cm]{geometry}
\usepackage{fancyhdr}
\usepackage{enumerate}
\usepackage[shortlabels]{enumitem}
\usepackage{parskip}
\usepackage[most]{tcolorbox}
\usepackage{float}
\usepackage{booktabs}
\usepackage{mathrsfs}
\usepackage{tikz}
\usepackage[hidelinks]{hyperref}


% cabecera
\pagestyle{fancy}
\fancyhead[l]{Mecánica de Medios Continuos}
\fancyhead[c]{Taller 4}
\fancyhead[r]{\today}
\fancyfoot[c]{\thepage}
\renewcommand{\headrulewidth}{0.2pt} % linea horizontal

\begin{document}
\tableofcontents

\section{Problema 1}

\section{Corrección del Quiz}

\section{Punto 1: Verdadero o Falso}

\subsection{Literal (a)}

\textbf{FALSO.} La distribución de velocidades moleculares en un gas clásico en equilibrio térmico es la distribución de Maxwell--Boltzmann:
$$
f(v)\,dv = 4\pi v^{2}\left(\frac{m}{2\pi k_{B}T}\right)^{3/2}\exp\!\left(-\frac{mv^{2}}{2k_{B}T}\right)dv,\qquad v\ge 0.
$$
La distribución de Poisson describe, en cambio, el número de moléculas contenidas en un volumen pequeño:
$$
\Pr(n\mid N)=\frac{N^{n}}{n!}\,e^{-N}.
$$
Son distribuciones con roles distintos: Maxwell--Boltzmann es una distribución continua de velocidades, mientras que Poisson es una distribución discreta de conteo.

\subsection{Literal (b)}

\textbf{VERDADERO.} La longitud de separación molecular es
$$
L_{\mathrm{mol}}=\left(\frac{V}{N}\right)^{1/3}=\left(\frac{M_{\mathrm{mol}}}{\rho N_{A}}\right)^{1/3}.
$$
En un líquido las moléculas están en contacto directo, de modo que $L_{\mathrm{mol}}$ es del orden del tamaño molecular; por ejemplo, para el agua líquida $L_{\mathrm{mol}}\approx 0{,}31\ \mathrm{nm}$. En un gas ideal, usando $pV=Nk_{B}T$,
$$
L_{\mathrm{mol}}=\left(\frac{k_{B}T}{p}\right)^{1/3},
$$
lo que a temperatura y presión ambiente da $L_{\mathrm{mol}}\approx 3{,}4\ \mathrm{nm}$, aproximadamente una décima parte mayor que en el líquido. Por tanto, en general $L_{\mathrm{mol}}^{\mathrm{gas}}>L_{\mathrm{mol}}^{\mathrm{líquido}}$.

\subsection{Literal (c)}

\textbf{FALSO.} La presión absoluta en un líquido incompresible rodeado de aire es
$$
p_{\mathrm{abs}}(z)=p_{\mathrm{atm}}+\rho g z.
$$
Si se duplica la profundidad $h$:
$$
p_{\mathrm{abs}}(2h)=p_{\mathrm{atm}}+2\rho g h.
$$
En cambio, el doble de la presión a profundidad $h$ sería
$$
2\,p_{\mathrm{abs}}(h)=2p_{\mathrm{atm}}+2\rho g h.
$$
Ambas expresiones difieren en $p_{\mathrm{atm}}$, por lo que $p_{\mathrm{abs}}(2h)\ne 2p_{\mathrm{abs}}(h)$. Lo que se duplica al duplicar la profundidad es la presión manométrica $\rho g h$, no la absoluta.

\section{Punto 2: Ecuación de estado del agua}

La ecuación de estado dada es
$$
p=p_{0}+B\left[\left(\frac{\rho}{\rho_{0}}\right)^{n}-1\right],
$$
con $B=3000\ \mathrm{atm}$, $n=7$, $p_{0}=1\ \mathrm{atm}$ y $\rho_{0}=1\ \mathrm{g/cm^{3}}$.

\subsection{Literal (a): Módulo de compresibilidad}

El módulo de compresibilidad se define como
$$
K=\rho\,\frac{dp}{d\rho}.
$$
Derivando la ecuación de estado:
$$
\frac{dp}{d\rho}=B\,n\left(\frac{\rho}{\rho_{0}}\right)^{n-1}\frac{1}{\rho_{0}}.
$$
Por tanto,
$$
K=nB\left(\frac{\rho}{\rho_{0}}\right)^{n}.
$$
En la superficie, donde $\rho=\rho_{0}$:
$$
K_{0}=nB=7\times 3000\ \mathrm{atm}=21\,000\ \mathrm{atm}\approx 2{,}1\times 10^{9}\ \mathrm{Pa}.
$$
Este valor es consistente con el módulo de compresibilidad conocido del agua ($\sim 2{,}1\ \mathrm{GPa}$).

\subsection{Literal (b): Perfiles de densidad y presión}

Se define la variable adimensional
$$
u(z)\equiv \frac{\rho(z)}{\rho_{0}},\qquad u(0)=1.
$$
La ecuación de estado se escribe
$$
p=p_{0}+B(u^{n}-1),
$$
de modo que
$$
\frac{dp}{dz}=Bn\,u^{\,n-1}\frac{du}{dz}.
$$
El equilibrio hidrostático exige
$$
\frac{dp}{dz}=-\rho g=-\rho_{0}g\,u.
$$
Igualando ambas expresiones:
$$
Bn\,u^{\,n-1}\frac{du}{dz}=-\rho_{0}g\,u
\quad\Longrightarrow\quad
u^{\,n-2}\,du=-\frac{\rho_{0}g}{Bn}\,dz.
$$
Para $n=7$ queda $u^{5}\,du=-\dfrac{\rho_{0}g}{7B}\,dz$. Integrando:
$$
\frac{u^{6}}{6}=-\frac{\rho_{0}g}{7B}\,z+C.
$$
Imponiendo $u(0)=1$ se obtiene $C=\tfrac{1}{6}$, y por tanto
$$
u^{6}=1-\frac{6\rho_{0}g}{7B}\,z.
$$
Los perfiles pedidos son entonces
$$
\rho(z)=\rho_{0}\left(1-\frac{6\rho_{0}g}{7B}\,z\right)^{1/6},
$$
$$
p(z)=p_{0}+B\left[\left(1-\frac{6\rho_{0}g}{7B}\,z\right)^{7/6}-1\right].
$$
El perfil existe mientras el argumento sea positivo, es decir, para
$$
0\le z<\frac{7B}{6\rho_{0}g}.
$$
En el límite $z\to 7B/(6\rho_{0}g)$, tanto $\rho$ como $p$ tienden a $0$ y $p_{0}$ respectivamente de manera consistente con la ecuación de estado.

\subsection{Literal (c): Bonificación --- Altura característica}

Una altura característica natural del problema es
$$
h_{c}=\frac{7B}{6\rho_{0}g}.
$$
Verificación dimensional: $[B]=\mathrm{Pa}=\mathrm{kg\,m^{-1}\,s^{-2}}$ y $[\rho_{0}g]=\mathrm{kg\,m^{-3}\,m\,s^{-2}}=\mathrm{Pa/m}$, de modo que
$$
[h_{c}]=\frac{\mathrm{Pa}}{\mathrm{Pa/m}}=\mathrm{m}.
$$
La elección es dimensionalmente consistente. Numéricamente:
$$
h_{c}=\frac{7\times 3000\times 101\,325\ \mathrm{Pa}}{6\times 1000\ \mathrm{kg/m^{3}}\times 9{,}81\ \mathrm{m/s^{2}}}
\approx 3{,}6\times 10^{4}\ \mathrm{m}\approx 36\ \mathrm{km}.
$$
Esta altura es mucho mayor que la profundidad real de los océanos ($\sim 11\ \mathrm{km}$), lo que justifica que el agua pueda tratarse como casi incompresible en la mayoría de aplicaciones oceánicas ordinarias. También puede compararse con la altura característica del caso estrictamente incompresible,
$$
h_{0}=\frac{p_{0}}{\rho_{0}g}\approx 10\ \mathrm{m},
$$
observando que $h_{c}=\dfrac{7B}{6p_{0}}\,h_{0}\approx 3500\,h_{0}$.

\end{document}
